\chapter{Bile Acids}

\section*{What Is This Test?}
Bile-acid testing measures circulating bile acids produced from cholesterol and used for digestion and absorption of fats. Serum bile acids can reflect impaired bile formation, bile flow, or hepatic clearance.

\section*{Alternative Names}
Serum Bile Acids; Total Bile Acids; Total Serum Bile Acids.

\section*{Why This Test Is Ordered}
Testing may support evaluation of cholestatic liver disease, unexplained itching or abnormal liver tests, and suspected intrahepatic cholestasis of pregnancy. It is interpreted with bile pigments, aminotransferases, alkaline phosphatase, clinical symptoms, and—when indicated—additional liver evaluation.

\section*{How This Test Is Performed}
A blood sample is collected, often after an overnight fast according to local laboratory instructions. The laboratory measures total or individual bile acids by an enzymatic or mass-spectrometric method.

\section*{Preparation, Risks, and Limitations}
Follow the laboratory's fasting and timing instructions. Venipuncture is the only usual risk. Meals, pregnancy, liver disease, medicines, and the timing of collection can affect results. A single result should not be interpreted without the gestational age or liver context when relevant.

\section*{Understanding Results}
Elevated bile acids suggest impaired hepatic handling or cholestasis but are not specific for one disease. In pregnancy, the concentration and trend help assess suspected cholestasis and guide specialist management; local pregnancy-specific thresholds should be used.

\section*{References}
Royal College of Obstetricians and Gynaecologists guidance on intrahepatic cholestasis of pregnancy.

\clearpage
\section*{Understanding Results and Follow-up}
The laboratory report should identify the specimen, method, units, reference interval or decision limit, and whether the result is positive, negative, abnormal, or inconclusive. Reference intervals vary with age, sex, pregnancy, specimen type, assay, and laboratory; a result outside a range is not automatically a diagnosis, and a result within a range does not always exclude disease. Discuss unexpected results with the ordering clinician, who can decide whether repeat or confirmatory testing, treatment, or referral is appropriate. Do not change medicines or delay urgent care based on an isolated result.


\section*{Regulatory Status and Authorization Date}
Regulatory status applies to the specific assay, instrument, manufacturer, and intended use rather than to the test name alone. No single approval date can be assigned to this test category. For a commercial assay, verify the current product in the relevant regulator's database: the U.S. Food and Drug Administration (FDA) clearance, approval, or authorization; India's Central Drugs Standard Control Organisation (CDSCO) licence or permission; the European Union In Vitro Diagnostic Regulation (EU IVDR) CE certificate issued through the applicable conformity-assessment route; or the United Kingdom Medicines and Healthcare products Regulatory Agency (MHRA) registration and UKCA/CE status. Laboratory-developed tests may not have an individual FDA, CDSCO, EU IVDR, or MHRA approval date.
\clearpage
